% GENERATED FILE. Edit canonical lessons, then run npm run generate:books.
% canonical-source-hash: fnv1a64:510dcb6723ad5eb5
% canonical-lessons: ES-C316-llanto, ES-C316-nunca, ES-C316-siempre, ES-C316-todo, ES-C316-nada

\chapter{Never, Always, All, Nothing}
\label{ch:nunca-y-siempre}
\hlchaptermodality{316}

\begin{chapteropening}
\textbf{By the end of this chapter:} \emph{I can answer with never, always, all and nothing in Spanish.}

Complete ES-C316-nada: put nunca and siempre against la risa and el llanto, then close the line with todo and nada.
\end{chapteropening}

\section[el llanto]{el llanto — the beating of the chest}

\label{lesson:ES-C316-llanto}



\begin{quote}
Which English words come from \emph{ridere}? (\emph{Ridicule}, \emph{deride}, \emph{risible}.) And what was a \emph{meticulous} Roman? (A frightened one.)
\end{quote}

\subsection*{What to know first}
\begin{itemize}
\raggedright
  \item la risa — the other half of this pair.
  \item el pecho — what the Latin says you are striking.
\end{itemize}

\begin{sounds}
\begin{itemize}
\raggedright
  \item \emph{LLAN-to}, two beats. The \emph{ll} is a \emph{y}, as in \emph{calle}. $\to$ reference
\end{itemize}
\end{sounds}

\begin{cousinweb}
\textbf{el llanto} — \textquotedblleft{}the weeping.\textquotedblright{}

Latin \textbf{plangere} meant \textquotedblleft{}to strike, to beat.\textquotedblright{} \textbf{Planctus} was the noun: a beating. And Roman mourning was a physical performance — you struck your own chest, hard, in public, where people could see you do it. The word for the grief is the word for the blow.

So \emph{el llanto} is built on \emph{el pecho}, which you learned five lessons ago. The chest is inside the etymology.

Two things happened on the way into Spanish. The ending wore off, and the \emph{pl-} turned into \emph{ll-}. That second change is a Spanish specialty and it is very regular. Latin \emph{pluvia} is \emph{la lluvia}. Latin \emph{plenus} is \emph{lleno}. Latin \emph{plorare} is \emph{llorar}. Latin \emph{flamma} is \emph{la llama}. Any Spanish word beginning \emph{ll-} is worth testing for a Latin \emph{pl-} or \emph{fl-} underneath it.

English kept \emph{plangere} twice. To \textbf{complain} is \emph{com-plangere}, to beat the breast thoroughly — which tells you how seriously the Romans took complaining. And \textbf{plangent}, said of a sound, means it strikes: a bell, a wave, something that hits and rings after.

Now set the pair side by side. \emph{La risa} and \emph{el llanto} are the two things a face does when it stops being neutral, and Spanish names them as a pair constantly.
\end{cousinweb}

\begin{grammarlens}[title={What you've built}]
The other half of the face, and the first word of a new chapter.

What follows is a set of very small words — the ones that say how often, how much, and whether at all. They are what you reach for when you answer. And the first two of them are exactly what anybody would say about laughing and weeping.
\end{grammarlens}

\subsection*{Your turn}
Say these aloud:

\begin{itemize}
\raggedright
  \item \emph{el llanto}
  \item \emph{la risa y el llanto}
  \item \emph{el miedo}, then \emph{el llanto}
\end{itemize}

\begin{tcolorbox}[breakable,colback=teal!4,colframe=teal!35!black,title={Before you move on}]
What did \emph{plangere} mean? (To strike.) What were Roman mourners beating? (Their own chest.) And what does a Spanish \emph{ll-} usually sit on in Latin? (A \emph{pl-} or an \emph{fl-}.)
\end{tcolorbox}
\section[nunca]{nunca — not ever, in one word}

\label{lesson:ES-C316-nunca}



\begin{quote}
What is a \emph{planctus}? (A beating.) And which English words come from \emph{ridere}? (\emph{Ridicule}, \emph{deride}.)
\end{quote}

\subsection*{What to know first}
\begin{itemize}
\raggedright
  \item no — the \emph{n} at the front of this one.
  \item el llanto — one of the things you can say this about.
\end{itemize}

\begin{sounds}
\begin{itemize}
\raggedright
  \item \emph{NUN-ca}, two beats. The \emph{c} is hard, as in \emph{cat}. $\to$ reference
\end{itemize}
\end{sounds}

\begin{cousinweb}
\textbf{nunca} — \textquotedblleft{}never.\textquotedblright{}

Latin \textbf{numquam} is two words fused together: \emph{ne} \textquotedblleft{}not\textquotedblright{} and \emph{umquam} \textquotedblleft{}ever.\textquotedblright{} Not-ever.

English \textbf{never} is also two words fused together: Old English \emph{ne} \textquotedblleft{}not\textquotedblright{} and \emph{æfre} \textquotedblleft{}ever.\textquotedblright{} Not-ever.

These two are not related. Latin and Germanic built them separately, out of their own parts, in the same order, to the same shape. When two languages reach the same construction by the same route without either copying the other, it usually means the route was the obvious one — and \textquotedblleft{}not ever\textquotedblright{} is about as obvious as an idea gets.

\emph{Nunca} does not change. No \emph{-o}, no \emph{-a}, nothing to agree with. It sits in front of a verb and negates it, and that is the whole of its job.
\end{cousinweb}

\begin{grammarlens}[title={What you've built}]
The first of the measuring words.

\emph{Nunca} is one end of a line. The next word is the other end.
\end{grammarlens}

\subsection*{Your turn}
Say these aloud:

\begin{itemize}
\raggedright
  \item \emph{nunca}
  \item \emph{no, nunca}
  \item \emph{el llanto}, then \emph{nunca}
\end{itemize}

\begin{tcolorbox}[breakable,colback=teal!4,colframe=teal!35!black,title={Before you move on}]
What two Latin words are inside \emph{numquam}? (\emph{Ne} and \emph{umquam} — not ever.) What two words are inside \emph{never}? (\emph{Ne} and \emph{ever}.) And does \emph{nunca} change for a woman? (No — it does not change at all.)
\end{tcolorbox}
\section[siempre]{siempre — the other end of the line}

\label{lesson:ES-C316-siempre}



\begin{quote}
What two pieces are inside \emph{numquam}? (\emph{Ne} and \emph{umquam}.) And what were Roman mourners striking? (Their chest.)
\end{quote}

\subsection*{What to know first}
\begin{itemize}
\raggedright
  \item nunca — the opposite end.
  \item la risa — something to measure with these two.
\end{itemize}

\begin{sounds}
\begin{itemize}
\raggedright
  \item \emph{SIEM-pre}, two beats. The \emph{ie} slides, and \emph{pr} runs together with no vowel between. $\to$ reference
\end{itemize}
\end{sounds}

\begin{cousinweb}
\textbf{siempre} — \textquotedblleft{}always.\textquotedblright{}

Latin \textbf{semper}, and this time almost nothing happened to it. The stressed \emph{e} broke to \emph{ie}, as it always does, and the ending softened. The meaning did not move at all.

English has no plain everyday word from it, but it has two showy ones. \textbf{Sempiternal} is \emph{semper} plus \emph{aeternus} — everlasting plus eternal, the same idea said twice, which is what a language does when one word is not emphatic enough. And \textbf{semper fidelis}, \textquotedblleft{}always faithful,\textquotedblright{} is the United States Marine Corps motto; most people who say \emph{semper fi} have never once thought about the first word.

\emph{Siempre}, like \emph{nunca}, refuses to change. \emph{Siempre} ends in \emph{-e}, which is the reason \emph{triste} already gave you. \emph{Nunca} ends in \emph{-a} and still does not move — because it is not describing a person at all. It is describing \emph{when}, and \emph{when} has no gender.
\end{cousinweb}

\begin{grammarlens}[title={What you've built}]
Both ends of the line.

\emph{Nunca} and \emph{siempre} between them cover the whole range. What comes next starts dividing up the middle.
\end{grammarlens}

\subsection*{Your turn}
Say these aloud:

\begin{itemize}
\raggedright
  \item \emph{siempre}
  \item \emph{sí, siempre}
  \item \emph{nunca}, then \emph{siempre}
\end{itemize}

\begin{tcolorbox}[breakable,colback=teal!4,colframe=teal!35!black,title={Before you move on}]
What is \emph{sempiternal} made of? (\emph{Semper} and \emph{aeternus} — everlasting twice.) What does \emph{semper fi} mean? (Always faithful.) And why does \emph{nunca} not change, even ending in \emph{-a}? (It describes when, not a person.)
\end{tcolorbox}
\section[todo]{todo — all of it, and this one does change}

\label{lesson:ES-C316-todo}



\begin{quote}
What is the first word of the Marine Corps motto? (\emph{Semper} — always.) And what two pieces are in \emph{never}? (\emph{Ne} and \emph{ever}.)
\end{quote}

\subsection*{What to know first}
\begin{itemize}
\raggedright
  \item mucho — a smaller helping than this.
  \item siempre — the same idea, applied to time.
\end{itemize}

\begin{sounds}
\begin{itemize}
\raggedright
  \item \emph{TO-do}, two beats. The \emph{d} between vowels is soft. $\to$ reference
\end{itemize}
\end{sounds}

\begin{cousinweb}
\textbf{todo} — \textquotedblleft{}all,\textquotedblright{} \textquotedblleft{}everything.\textquotedblright{}

Latin \textbf{totus} meant whole, entire. English has it undisguised: a \textbf{total} is the whole of something, and \textbf{totally} means wholly. There is no sound change to unpick here. The word walked into Spanish and dropped one \emph{t} on the way.

Portuguese took the same Latin word down its own road and arrived at \emph{tudo}. Spanish \emph{todo} and Portuguese \emph{tudo} are siblings, not borrowings — each is \emph{totus}, worn down by its own language's habits.

Now the useful part. \emph{Nunca} and \emph{siempre} do not change. \textbf{Todo does.}

\emph{Todo} is describing a thing, and in Spanish anything that describes a thing has to agree with it. \emph{Todo} in front of a masculine word; \emph{toda la noche} — all the night — in front of a feminine one. The \emph{-o} becomes \emph{-a} because \emph{noche} is feminine.

That is the rule \emph{triste} held out on. Words ending in \emph{-o} switch to \emph{-a}. Words ending in \emph{-e} do not move. \emph{Todo} ends in \emph{-o}, so it switches.
\end{cousinweb}

\begin{grammarlens}[title={What you've built}]
The first word in this chapter that agrees.

\emph{Nunca} and \emph{siempre} measure time and stay put. \emph{Todo} measures a thing, so it has to match the thing.
\end{grammarlens}

\subsection*{Your turn}
Say these aloud:

\begin{itemize}
\raggedright
  \item \emph{todo}
  \item \emph{toda la noche}
  \item \emph{todo}, then \emph{toda}
\end{itemize}

\begin{tcolorbox}[breakable,colback=teal!4,colframe=teal!35!black,title={Before you move on}]
What did \emph{totus} mean? (Whole.) Which English word is it, barely changed? (\emph{Total}.) And what does \emph{todo} become in front of \emph{la noche}? (\emph{Toda}.)
\end{tcolorbox}
\section[nada]{nada — a thing born, worn down to nothing}

\label{lesson:ES-C316-nada}



\begin{quote}
What does \emph{todo} become before \emph{la noche}? (\emph{Toda}.) And what is \emph{sempiternal} made of? (\emph{Semper} and \emph{aeternus}.)
\end{quote}

\subsection*{What to know first}
\begin{itemize}
\raggedright
  \item de nada — you have been saying this word since chapter four.
  \item todo — its exact opposite.
\end{itemize}

\begin{sounds}
\begin{itemize}
\raggedright
  \item \emph{NA-da}, two beats. Both \emph{a}s are open and the \emph{d} is soft. $\to$ reference
\end{itemize}
\end{sounds}

\begin{cousinweb}
\textbf{nada} — \textquotedblleft{}nothing.\textquotedblright{}

This is the best story in the chapter.

Latin had a phrase: \textbf{res nata}, \textquotedblleft{}a thing born.\textquotedblright{} To say there was nothing at all, you said \emph{non res nata} — not a thing born, not one single thing that has come into being. Perfectly ordinary emphatic negation, the kind English still does with \textquotedblleft{}not a living soul.\textquotedblright{}

Then the phrase wore down. \emph{Res} fell off the middle. \emph{Non} drifted away and became the separate \emph{no}. What was left standing was \textbf{nata} — the participle meaning \emph{born} — carrying, entirely on its own, the meaning \textquotedblleft{}nothing.\textquotedblright{}

So the Spanish word for nothing is, historically, the word for a thing that exists. It means its own opposite, because the negative that made it negative has fallen off the front of it.

\emph{Nata} is the feminine of \emph{natus}, born, and that root is everywhere in English: \textbf{natal}, \textbf{native}, \textbf{nativity}, \textbf{nature}, \textbf{nascent}, \textbf{innate}, \textbf{renaissance}. Every one of them is about coming into being. And \emph{nada} is the same word.

You have been saying it since chapter four. \emph{De nada} — \textquotedblleft{}it's nothing\textquotedblright{} — is \emph{of nothing}, and now you know what the nothing was made of.

Keep one thing for the next chapter: \emph{nada} has a twin, and the twin is about people.
\end{cousinweb}

\begin{grammarlens}[title={What you've built}]
Fifteen words, and this chapter closes on its best etymology.

\emph{Todo} and \emph{nada} are the two ends of it. Everything, and not a thing born.
\end{grammarlens}

\subsection*{Your turn}
Say these aloud:

\begin{itemize}
\raggedright
  \item \emph{nada}
  \item \emph{de nada}
  \item \emph{todo}, then \emph{nada}
\end{itemize}

\begin{tcolorbox}[breakable,colback=teal!4,colframe=teal!35!black,title={Before you move on}]
What was \emph{res nata}? (A thing born.) Which part of the phrase survived as \emph{nada}? (\emph{Nata} — the word for born.) And name an English word from the same root. (\emph{Native}, \emph{nature}, \emph{natal} — any of them.)
\end{tcolorbox}
